\section{Limitations}

The current strongest performance evidence is in-domain ICLR 2024.
The ICLR 2025 repro set is validated but small, and expanded80 plus NeurIPS 2024 limit100 are active frontiers rather than random or natural-prevalence samples.
The ICLR 2023 random80 slice adds random/stratified evidence within the ICLR venue family, but it is still standard single-user validation and not a full natural-prevalence estimate for all venues.
Most labels are standard single-user confirmed labels.
Second-pass reliability evidence is currently bounded to two packets: an independent mini60 packet and a broader boundary160 packet that was prelabel-assisted before user confirmation.
We therefore frame cross-year and cross-venue evidence as bounded transfer and external-validity evidence, not as a complete estimate of performance on all issues.
For the same reason, frontier failure-taxonomy counts should be read as diagnostic counts on intentionally hard frontiers, while random80 counts should be read by measured slice design rather than as broad prevalence estimates.
These agreement results should be interpreted narrowly: they are bounded consistency checks on hard cases, not a substitute for full two-annotator coverage across all packets.

The label rubric is issue-level and evidence-based, but some concerns remain ambiguous.
For example, a narrowed claim may make an original criticism less severe without fully resolving it.
These cases depend on careful evidence spans and notes.

The current labeled pool has too few regressed examples for stable regression-performance conclusions.
We include the label because it is conceptually important, but treat regression performance as provisional until more examples are collected.

RevTrack is also text-first.
The current pipeline uses review text, author responses, revision summaries, and available paper metadata rather than full PDF differencing.
This keeps the first benchmark auditable and lightweight, but it can miss changes that are only visible in figures, equations, or revised appendices.
Future versions should add PDF-aware evidence extraction while preserving issue-level provenance.

Finally, the current standard labels come from user-reviewed AI-assisted signoff and user-confirmed expanded80, NeurIPS 2024, and ICLR 2023 adjudication.
A separate independent mini60 second pass is complete, and a user-confirmed boundary160 second-pass packet is also complete.
Because boundary160 was prelabel-assisted, full-scale blind-independent two-annotator coverage remains future work.
